\section{The Event of 19 September 1846 and Its First Transmission}

\subsection{Why the transmission has to come first}

An apparition study is only as good as the chain by which the
reported words reached the page. At La Salette that chain is short
in one respect and awkward in another. Short: the two children told
what they had seen on the evening of the same day, in the kitchen of
the farm where one of them was in service, and the parish priest
announced it from the pulpit the next morning. Awkward: the children
were illiterate, spoke an Alpine \term{patois} and very little
French, and everything now read as ``the message'' passed through
adults who wrote in French what children said partly in a language
those adults had to translate.

This study therefore names its strata before it narrates, and keeps
them named throughout:

\begin{itemize}[leftmargin=1.2em,itemsep=0.15em,topsep=0.25em]
\item \textbf{Same-day oral transmission (19--20 September 1846).}
The children's recital to their employers, to the parish priest of
La Salette, and to the mayor. No autograph survives from the
children, who could not write.
\item \textbf{Early adult records (from late 1846).} Written
relations by local priests and laymen, the interrogations conducted
by the archpriest of Corps and others, and the notes associated with
Abbé Lagier. This project did not reach these documents in their
originals or in a critical edition; it reached them only where a
later witness reproduces them.
\item \textbf{The diocesan inquiry (1847).} Depositions signed 27
November 1847 by the children's employers, the commission sittings
of November--December 1847 held in the bishop's presence, and the
\term{rapporteur}'s printed report.
\item \textbf{Period apologetic literature (1850s).} Books that
reproduce the foregoing in translation and with argument---the class
to which the two 1854 witnesses most used here belong.
\item \textbf{Received custodial text (present).} The French
discourse as the Sanctuary and the Missionaries now publish it, a
harmonized text without an apparatus.
\end{itemize}

Nothing in this study treats the fifth stratum as though it were the
first. Where the received text says more than the earliest reachable
record, \S3 says so.

\subsection{Place, persons, and the days before}

The site is a high pasture above the village of La Salette, in the
canton and archpresbytery of Corps, in the Alps of the Isère, at the
head of a ravine where a small intermittent spring ran. The parish
belonged then, as the whole dossier does now, to the diocese of
Grenoble, today Grenoble--Vienne.

Maximin Giraud was born at Corps on 26 August 1835, son of a
wheelwright; his mother died when he was seventeen months old and he
was brought up, on the Sanctuary's own frank account, ``like a wild
plant.'' He was eleven. Mélanie Calvat, called Mathieu, was born
7 November 1831, the fourth of ten children of a sawyer; she had
been put out to service to mind cows from the age of ten, at Quet
and at Sainte-Luce, before reaching the farm of Jean-Baptiste Pra at
the hamlet of Les Ablandins. She was going on fifteen. Both families
lived at Corps, at opposite ends of the town; the children had not
previously known each other. Maximin had come up only days before,
to replace for a week the sick herdsman of Pierre Selme, Pra's
neighbour.\endnote{Sanctuaire Notre-Dame de La Salette, ``Les
témoins: Mélanie et Maximin,'' the Sanctuary's own biographical
page, read 2026-07-25 in the Internet Archive capture of 8 April
2013 (References); the Sanctuary's current site carries the event
and recognition pages but no longer this biography, and the capture
is cited with its date. The service history and the children's
non-acquaintance are also in the 1847 depositions and in the
\term{rapporteur}'s note from Abbé Lagier's papers, as reproduced in
Ullathorne 1854, pp.~30--32.}

The children could barely speak French. An English bishop who
travelled to Corps in 1854 and read the diocesan documents put it
bluntly: ``these children neither spoke nor understood much of the
French language. Their own is a singular patois\ldots\ Thus, they
knew not the French for so common an object as the
potato.''\endnote{Ullathorne 1854, p.~32. The observation matters
for \S3: the received discourse contains a passage in which the
Lady, having spoken French, repeats herself in \term{patois}
precisely because the children have not understood.}

\subsection{The encounter, as the earliest reachable recital gives
it}

The sequence, in the recital the children gave and kept giving: they
had brought their cows to the high pasture, eaten, and slept a
little; Mélanie woke first, did not see the cows, woke Maximin, and
they crossed the little stream and saw the animals lying on the far
slope; turning back for their bags near the dried-up spring, Mélanie
saw a brightness ``like the sun, it was far more brilliant, but it
had not the same colour.'' In the light was a Lady seated, her head
in her hands. Mélanie let her stick fall in fear; Maximin said he
would strike whatever it was if it did anything. The Lady rose,
crossed her arms, and spoke: ``Come near, my children, be not
afraid, I am here to tell you great news.''\endnote{Mélanie's
recital, Ullathorne 1854, pp.~35--36, printing the version the
witness says the children ``gave\ldots\ on the same evening to their
employers; as they gave it the next morning to the Curé of La
Salette; as Mélanie recounted it the same day to Monsieur Peytard,
the mayor of La Salette,'' and afterwards constantly. That claim of
constancy is the witness's own and is treated here as his claim, not
as established fact; what the deposition-level evidence establishes
is examined at \S2.4.}

The discourse followed (\S3). At its end the Lady told them twice,
in French, to make it known to all her people; she crossed the
stream, climbed the slope without bending the grass, rose a little
above the ground, looked to heaven and then to earth, and
disappeared into the light. The children's reported exchange
afterwards is the sort of detail no forger invents: Mélanie said
``perhaps it is a great saint''; Maximin answered that if they had
known it was a great saint they would have asked her to take them
with her; he then reached out to catch some of the remaining
brightness, and there was nothing there.\endnote{Ullathorne 1854,
pp.~42--43.} The figure's described appearance---white robe, a cape
with roses, a crown of roses, gold-coloured apron, and on the breast
a crucifix with a hammer on one side and pincers on the other, hung
from a small chain, with a larger chain falling from the ends of the
cross---is the description Mélanie gave under interrogation, and is
the source of every later image and of the Sanctuary's present
account of the light as coming from the crucifix
itself.\endnote{Mélanie's answer under interrogation, Ullathorne
1854, p.~44 (``She had white shoes, with roses round them of all
colours\ldots\ a crown with roses round her cap. She wore a very
small chain, on which was hung a cross, with a figure of our Lord;
on the right were pincers, on the left a hammer\ldots''). The
Sanctuary's present formulation---``Toute la clarté dont elle est
formée et qui les enveloppe tous les trois, vient d'un grand
Crucifix qu'elle porte sur sa poitrine, entouré d'un marteau et de
tenailles''---is a custodial development of the same description and
is quoted as such
(\url{https://lasalette.cef.fr/lhistoire/}).}

\subsection{The evening at Les Ablandins, and the depositions of 27
November 1847}

The best evidence about the first hours is not the children's own
narrative but the sworn statements of the two farmers, signed 27
November 1847 for the diocesan inquiry. Pierre Selme deposed that
Maximin failed to return to him in the field, that at evening the
boy said ``you don't know what has happened,'' recited what he had
seen, and that the next morning the two employers sent both children
to the curé of La Salette, ``who, the same day, at Mass, made known
to his parishioners what they had seen and heard.'' Selme added the
observation that matters most evidentially: during the four and a
half days the boy guarded his cows, in open unfenced pasture where a
standing man sees the whole field, ``I never saw either priest or
layman come near him or speak to him,'' and Mélanie in the adjoining
field was likewise always alone.

\begin{witnesstext}{Jean-Baptiste Pra, deposition signed 27 November
1847 (1854 English witness)}
``On Saturday, the 19th Sept., 1846, they both came and told me what
they had seen and heard on the level. Next morning they went
together to the Curé of La Salette, who the same day at Mass told it
from the pulpit to his parishioners\ldots\ I ought to add, before
signing, that for some days at first I did not believe the statement
of the children, and I several times tried to induce little Mélanie
to take the money offered her to keep silence. The child constantly
refused the money offered to her. She has always persisted, in spite
of threats as well as promises of reward. The mayor of La Salette,
amongst others, tried all sorts of means in vain to put the child in
contradiction with herself; he could not succeed. He then offered
her money; she refused it, and she replied to his menaces, that she
would always and everywhere repeat what the Blessed Virgin had said
to her.''\endnote{Ullathorne 1854, pp.~30--31, printing the
depositions of Pierre Selme and Jean-Baptiste Pra ``signed the 27th
of November, 1847.'' This study quotes the 1854 English witness and
does not assert the French original's wording. Pra's date ``19
Sept., 1846'' is printed in the scan as ``1840''; the reading 1846
is required by the deposition's own context and by every other
witness, and the scan defect is recorded rather than silently
corrected.}
\end{witnesstext}

Two things follow. First, the earliest transmission is documented
and rapid: same-day telling, next-morning pulpit announcement, and
an unbelieving employer as a hostile-turned-corroborating witness.
Second, the resistance to money and threats is attested from outside
the children, by a man who says he himself tried the money.

\subsection{The children as the inquiry found them}

The commission of 1847 also took the opinion of the superioress of
the convent school at Corps, where both children had been placed
some months after the event. Of Maximin she said he had ``but
ordinary abilities,'' was ``sufficiently obedient, but light, fond of
play, and always in motion,'' and that ``he never speaks to us of
the affair of La Salette, and we avoid the subject with him, that he
may not take any importance to himself.'' The same institution's
early anxiety about Mélanie is preserved by the Sanctuary itself:
already in November 1847 her superior feared ``que Mélanie ne tirât
vanité de la position que l'événement lui a
faite.''\endnote{Superioress's opinion: Ullathorne 1854, pp.~32--33.
Mélanie's superior's 1847 anxiety: Sanctuary, ``Les témoins''
(archived capture, note~3), which reports it in the Sanctuary's own
voice and adds its own explanation---a poor girl deprived of
affection, in service from the age of ten, suddenly under public
attention. The Sanctuary's frankness about its own seers is a
notable feature of the custodial record and is used in \S7.}

That is an unusual pair of contemporary judgments to find in a
dossier that a diocese would later approve: the boy is scattered and
must be kept from self-importance; the girl is already, in the
judgment of the religious who taught her, at risk of vanity from
the event. The 1851 mandement did not judge the children's later
lives. It judged an apparition on a named day. \S7 shows how
completely those two objects came apart.

\subsection{What the event corpus of this study is}

The included corpus is exactly one dossier: the single reported
apparition of 19 September 1846 to Maximin Giraud and Mélanie
Calvat on the mountain above La Salette. There is no series, no
fortnight, no set of later appearances forming part of the judged
event. Later reported locutions, visions, and revelations attributed
to either seer---including everything in Mélanie's expanded
publications (\S5)---lie outside the 1851 act and are treated in
this study as documents about their authors, not as parts of the
event. Reported cures and prodigies at the spring are named by the
mandement as corroborating evidence in article~2 (``the multitude of
prodigies which have been the consequence of the said event''); this
study does not certify, enumerate, or re-examine any of them, and
records that it located no institutional list of individually
recognized La Salette miracles comparable to those published for
other shrines.\endnote{Bounded negative: no list of individually
recognized miraculous cures was found on the Sanctuary's current
site or on the Missionaries' French and international sites (pages
read 2026-07-25 and enumerated in \work{research/source-audit.md}).
The mandement's own article~2 language is from the Italian witness
(Rousselot 1854, p.~44) and Ullathorne 1854, p.~120. Absence of a
published list is not evidence that no diocesan act ever recognized
a cure; it is the boundary of this project's search.}
